% =============================================================================
% Section 4 — Evaluation
% =============================================================================

\section{Evaluation}
\label{sec:evaluation}

We evaluated NeuroAssist AI through a~preliminary case study involving one
advanced-stage MND patient and five caregivers and volunteers. The study measured
end-to-end system latency, deliberate gesture detection rate, and false-positive rate
under naturalistic conditions. Results should be interpreted as a~preliminary case
study, not a~clinical trial (see Section~\ref{sec:discussion}).

% ---------------------------------------------------------------------------
\subsection{Participants and Ethics}
\label{sec:participants}

Six participants took part in the study (Table~\ref{tab:participants}).
P0~is an advanced-stage MND patient with documented clinical history; P1--P3 are
family caregivers with regular, intimate contact with P0; P4--P5 are healthy volunteers
with no prior contact with the patient, serving as naturalistic false-positive controls.

\begin{table}[!t]
  \renewcommand{\arraystretch}{1.15}
  \caption{Study Participants}
  \label{tab:participants}
  \centering
  \begin{tabular}{llll}
    \toprule
    ID & Role & Relationship to P0 & Session \\
    \midrule
    P0 & MND patient          & ---                              & Calibration + detection \\
    P1 & Spouse / home carer  & Intimate, daily (physiotherapist) & FP + real alarm \\
    P2 & Caregiver 1          & NDIS professional, frequent      & FP + real alarm \\
    P3 & Caregiver 2          & NDIS professional, frequent      & FP + real alarm \\
    P4 & Control           & None                 & Naturalistic FP baseline \\
    P5 & Control           & None                 & Naturalistic FP baseline \\
    \bottomrule
  \end{tabular}
\end{table}

All participants provided informed consent prior to their session. The study was
conducted in P0's home in Melbourne, Australia. P1 is P0's spouse and a~trained
physiotherapist providing daily home care (ineligible for NDIS funding as a~family
member under Australian disability support policy); P2 and P3 are National Disability
Insurance Scheme (NDIS)-funded professional support workers with regular structured
contact with P0. The study constitutes a~pilot
evaluation; no formal human research ethics committee (HREC) approval was sought for
this preliminary case study scope. Full clinical evaluation with a~statistically representative sample of MND
patients will require HREC approval and is planned as future work.

P0 presents with truncal-onset MND (symptom onset: late 2019; disease duration at evaluation: approximately 6.5~years), four-limb weakness with preserved voluntary eyelid movement, and non-invasive ventilation (NIV) prescribed since early 2026 and active at the time of the study. Forced vital capacity was 71\,\% and sniff nasal inspiratory pressure (SNIP) 31\,\% at last documented pre-experiment respiratory assessment (May 2023), with slow disease progression confirmed at that time. \fillme{Confirm before submission: ALSFRS-R score if available; current palpebral mobility assessment.} This clinical profile informed the calibration protocol described in Section~\ref{sec:calibration}.

% ---------------------------------------------------------------------------
\subsection{Experimental Setup}
\label{sec:setup}

\textbf{Hardware.} Each session used a~standard USB webcam
(\fillme{model/resolution, e.g.\ Logitech C920, 1080p@30fps}) mounted
at eye level approximately 50--70\,cm from the participant's face, connected to a
laptop running the NeuroAssist AI stack locally. No external network access was
required during sessions.

\textbf{Software.} The edge detection pipeline ran as
\texttt{poc-edge-agent-webcam.py} (Python~3.11), communicating with the Angular
web application (\texttt{poc-gesture-detector.ts}) via \texttt{localhost}. The
\texttt{evnt-service} and \texttt{gest-service} microservices ran locally via Docker
Compose. Measurement scripts \texttt{measure\_latency.py} and
\texttt{false\_positive\_experiment.py} recorded all timestamps to CSV; alarm
delivery timestamps were captured by an Angular \texttt{effect()} hook writing to
\texttt{localStorage} on each alarm render.

\textbf{Environment.} Sessions were conducted in P0's home during
\fillme{time of day}. Ambient illumination was natural/artificial
indoor lighting; lux levels were not formally measured. Camera position was recorded
per session.

% ---------------------------------------------------------------------------
\subsection{Patient Calibration}
\label{sec:calibration}

Before formal trials, P0's \texttt{DetectionProfile} was calibrated by observing
natural palpebral movement for approximately 5--10~minutes. P0's clinical profile — progressive MND with documented respiratory muscle involvement (FVC 71\,\%, SNIP 31\,\%) and NIV prescribed since early 2026 — is representative of the population for whom NeuroAssist AI is principally designed~\cite{berlowitz2021niv, yang2026pva}, and the calibrated DetectionProfile reflects the residual palpebral mobility characteristic of this stage. The calibration procedure adjusted two parameters: $\theta_\text{close}$ (EAR threshold for
blink onset) and \texttt{BLINK\_MIN\_FRAMES} (minimum consecutive frames below
$\theta_\text{close}$ required to register a~blink).

Table~\ref{tab:calibration} reports the calibrated parameters for P0 alongside
the defaults used for P4 and P5.

\begin{table}[!t]
  \renewcommand{\arraystretch}{1.15}
  \caption{DetectionProfile Calibration Parameters}
  \label{tab:calibration}
  \centering
  \begin{tabular}{lccc}
    \toprule
    Participant & $\theta_\text{close}$ & $\theta_\text{open}$ & Min.\ frames \\
    \midrule
    P0 (MND patient) & \textbf{[X.XX]} & \textbf{[X.XX]} & \textbf{[N]} \\
    P4--P5 (controls, default) & 0.20 & 0.27 & 3 \\
    \bottomrule
  \end{tabular}
\end{table}

% ---------------------------------------------------------------------------
\subsection{Latency}
\label{sec:latency}

End-to-end latency was decomposed into two segments: (i)~gesture-to-HTTP, measured
as the elapsed time between the state machine entering \texttt{DONE} and the
\texttt{POST /api/v0/report-ingest} response being received by the edge agent; and
(ii)~event-to-alarm-visible, measured as the elapsed time between the
\texttt{createdAt} timestamp in the \texttt{evnt-service} response and the
\texttt{alarmShownAtMs} value recorded in \texttt{localStorage} by the Angular
application. Total end-to-end latency is the sum of these two segments plus any
HTTP round-trip overhead on the local network.

Table~\ref{tab:latency} reports latency statistics across all trials.

% TODO: replace this placeholder with \input after running generate_tables.py
% \input{../../../scripts/paper/data/table_latency}
\begin{table}[!t]
  \renewcommand{\arraystretch}{1.15}
  \caption{End-to-End Latency (ms)}
  \label{tab:latency}
  \centering
  \begin{tabular}{lcccc}
    \toprule
    Segment & $N$ & Mean & $\sigma$ & Median \\
    \midrule
    Gesture $\rightarrow$ HTTP  & [N] & [X] & [Y] & [Z] \\
    Event $\rightarrow$ alarm   & [N] & [X] & [Y] & [Z] \\
    \textbf{End-to-end total}   & [N] & \textbf{[X]} & [Y] & [Z] \\
    \bottomrule
  \end{tabular}
\end{table}

Mean end-to-end latency was \textbf{[X]}\,ms ($\sigma$\,=\,[Y]\,ms, median
[Z]\,ms, $N$\,=\,[N] trials). The event-to-alarm segment accounted for the
largest share of total latency, consistent with the 5-second polling interval
used in the PoC caregiver application (worst-case contribution: 5{,}000\,ms; see
Section~\ref{sec:discussion} for the planned WebSocket mitigation). The
gesture-to-HTTP segment was dominated by local HTTP overhead and remained below
[X]\,ms across all trials, confirming that the state machine and WASM pipeline
introduce negligible latency.

% ---------------------------------------------------------------------------
\subsection{False-Positive Rate}
\label{sec:fpr}

False-positive rate (FPR) was measured as the number of unintended alarm triggers
per minute of continuous facial activity, for both the six-state deliberate gesture
state machine (SM) and a~simple EAR threshold baseline (threshold $\leq
\theta_\text{close}$ for $\geq 1$ frame). Each session ran for 10\,minutes
($N$\,=\,600\,s) of naturalistic facial activity without the participant attempting
to trigger the gesture.

Table~\ref{tab:fp} reports FPR per participant and per group.

% TODO: replace with \input after running generate_tables.py
% \input{../../../scripts/paper/data/table_fp}
\begin{table}[!t]
  \renewcommand{\arraystretch}{1.15}
  \caption{False-Positive Rate (triggers/min)}
  \label{tab:fp}
  \centering
  \begin{tabular}{lccc}
    \toprule
    Participant & Group & Simple threshold & State machine \\
    \midrule
    P1 & Caregiver & [A] & [B] \\
    P2 & Caregiver & [A] & [B] \\
    P3 & Caregiver & [A] & [B] \\
    P4 & Control   & [A] & [B] \\
    P5 & Control   & [A] & [B] \\
    \midrule
    \textbf{Caregiver mean} & & [A] & \textbf{[B]} \\
    \textbf{Control mean}   & & [A] & \textbf{[B]} \\
    \bottomrule
  \end{tabular}
\end{table}

The six-state state machine reduced FPR from [A]\,triggers/min (simple threshold)
to [B]\,triggers/min averaged across all participants, a~relative reduction of
[C]\%. Participants P1--P3, who interacted with P0 in proximity contexts spanning intimate
spousal care (P1) and structured professional NDIS support (P2--P3), exhibited
\textbf{[higher/lower/comparable]} FPR compared to control participants (P4--P5)
engaged in neutral tasks. This three-tier contrast — intimate domestic care, professional
support care, and neutral baseline — provides a~novel ecological gradient for studying
FPR variation, not previously reported in the AAC false-positive literature.

% ---------------------------------------------------------------------------
\subsection{Detection Rate}
\label{sec:detection}

Detection rate (recall) was measured as the proportion of deliberate double-blink
gestures that successfully triggered the state machine and produced a~confirmed
alarm. P0 performed $N$\,=\,[N] deliberate double-blink attempts over the session
with rest pauses of at least 30\,s between attempts. Control participants (P4, P5)
performed the same protocol for comparison with a~healthy-volunteer baseline.

P0 achieved a~detection rate of \textbf{[D]\%} ([M] of [N] attempts detected),
using the calibrated \texttt{DetectionProfile} from Section~\ref{sec:calibration}.
Of the [N\,-\,M] missed attempts, [X] were attributed to incomplete eyelid closure
below the calibrated $\theta_\text{close}$ (observable from the EAR trace log),
and [Y] to \fillme{other cause, e.g.\ head movement causing landmark loss}. Control participants P4 and P5 achieved detection rates of [D1]\% and
[D2]\%, respectively, using default parameters. The higher detection rate of
healthy volunteers, if confirmed, is consistent with greater residual eyelid
mobility compared to advanced-stage MND — a~difference the patient-specific
calibration protocol aims to partially compensate.